% Source: https://artofproblemsolving.com/community/c3h2477986_how_to_do_reciprocals
% Stable identity: aops-topic-2477986-q01
\textbf{(a)} $\dfrac17$; $\dfrac53$ (since $\dfrac35\cdot\dfrac53=1$); $2\dfrac14=\dfrac94$, so its reciprocal is $\dfrac49$; and $0.4=\dfrac25$, so its reciprocal is $\dfrac52=2.5$.

\textbf{(b)} Dividing $a$ by $c$ means finding the number $q$ with $q\cdot c=a$. The number $q=a\cdot\dfrac1c$ works, because $a\cdot\dfrac1c\cdot c=a\cdot1=a$, and it is the only one: if $qc=a$ then $q=qc\cdot\dfrac1c=a\cdot\dfrac1c$. So $a\div c=a\cdot\dfrac1c$. Hence
\[\dfrac56\div\dfrac{10}9=\dfrac56\cdot\dfrac9{10}=\dfrac{45}{60}=\boxed{\dfrac34}.\]

\textbf{(c)} If $x=\dfrac1x$ then $x^2=1$, so $x=1$ or $x=-1$, and both work. ($0$ has no reciprocal.) The answer is $\boxed{x=\pm1}$.
